\documentclass[12pt,a4paper]{article}
\usepackage[utf8]{inputenc}
\usepackage[T1]{fontenc}
\usepackage[english]{babel}
\usepackage{amsmath}
\usepackage{amsfonts}
\usepackage{amssymb}
\usepackage{graphicx}
\usepackage{float}
\usepackage{listings}
\usepackage{xcolor}
\usepackage{hyperref}
\usepackage{geometry}
\geometry{margin=1in}

\lstset{
    language=C,
    basicstyle=\ttfamily\footnotesize,
    keywordstyle=\color{blue},
    commentstyle=\color{green!60!black},
    stringstyle=\color{red},
    numbers=left,
    numberstyle=\tiny,
    frame=single,
    breaklines=true,
    captionpos=b
}

\title{CSE344 System Programming Homework 5 Report}
\author{Student ID: 220104004043}
\date{\today}

\begin{document}

\maketitle

\section{Introduction}

This report presents the implementation of a cargo delivery simulator for CSE344 System Programming Homework 5. The simulator uses POSIX threads to manage multiple courier threads that process delivery orders from a priority queue. The system supports order prioritization, atomic statistics tracking, and safe shutdown via SIGINT signal handling.

\section{Problem Description}

The homework requires implementing a multi-threaded cargo delivery system with the following features:

\begin{itemize}
    \item Multiple courier threads processing orders from a shared priority queue
    \item Priority-based ordering (EXPRESS > STANDARD > ECONOMY, then by order ID)
    \item Atomic statistics tracking for completed/cancelled orders and delivery times
    \item SIGINT-safe shutdown with proper cancellation of pending orders
    \item Required logging format for all operations
    \item Memory safety and thread synchronization
\end{itemize}

\section{Design and Architecture}

\subsection{Thread Architecture}

The system consists of the following threads:

\begin{itemize}
    \item \textbf{Main Thread}: Initializes resources, loads orders, creates courier threads, waits for completion, and handles final statistics
    \item \textbf{Courier Threads}: Process orders from the priority queue, simulate deliveries, and update statistics
    \item \textbf{Signal Thread}: Dedicated thread for handling SIGINT signals safely
\end{itemize}

\subsection{Data Structures}

\subsubsection{Priority Queue}
A dynamic array-based priority queue that maintains sorted order by priority and ID:

\begin{lstlisting}
typedef struct {
    Order *orders;
    uint32_t count;
    uint32_t capacity;
} PriorityQueue;
\end{lstlisting}

\subsubsection{Order Structure}
Represents a delivery order with ID, recipient name, priority, and duration:

\begin{lstlisting}
typedef struct {
    uint32_t id;
    char name[33];
    Priority priority;
    uint32_t duration;
} Order;
\end{lstlisting}

\subsubsection{Global Statistics}
Atomic counters for thread-safe statistics tracking:

\begin{lstlisting}
typedef struct {
    _Atomic(uint32_t) completed_orders;
    _Atomic(uint32_t) cancelled_orders;
    _Atomic(uint64_t) total_delivery_time;
    CourierStats *couriers;
    uint32_t num_couriers;
    uint32_t total_orders;
} GlobalStats;
\end{lstlisting}

\subsection{Synchronization}

The system uses the following synchronization primitives:

\begin{itemize}
    \item \textbf{Mutexes}: Protect shared data structures (queue, log output)
    \item \textbf{Condition Variables}: Coordinate between producer/consumer threads
    \item \textbf{Atomic Variables}: Thread-safe counters without locks
\end{itemize}

\section{Implementation Details}

\subsection{Order Parsing}

The \texttt{order\_parse\_line} function validates and parses input lines with strict format checking:

\begin{itemize}
    \item ID: positive integer
    \item Name: ASCII alphanumeric + underscore, max 32 characters
    \item Priority: EXPRESS/STANDARD/ECONOMY
    \item Duration: positive integer
\end{itemize}

\subsection{Priority Queue Operations}

The queue maintains sorted order through insertion sort during enqueue operations:

\begin{lstlisting}
int pq_insert(PriorityQueue *pq, const Order *order) {
    // Find insertion position
    uint32_t pos = pq->count;
    for (uint32_t i = 0; i < pq->count; i++) {
        if (order_compare(order, &pq->orders[i]) < 0) {
            pos = i;
            break;
        }
    }
    // Shift and insert
    // ...
}
\end{lstlisting}

\subsection{Thread Management}

\subsubsection{Courier Thread Logic}

Each courier thread follows this loop:

\begin{enumerate}
    \item Wait for orders in queue (protected by mutex + condition)
    \item Dequeue highest priority order
    \item Simulate delivery using \texttt{nanosleep}
    \item Update statistics atomically
    \item Repeat until shutdown
\end{enumerate}

\subsubsection{Signal Handling}

A dedicated signal thread uses \texttt{sigtimedwait} to handle SIGINT:

\begin{lstlisting}
static void *signal_thread_func(void *arg) {
    sigset_t sigset;
    sigemptyset(&sigset);
    sigaddset(&sigset, SIGINT);
    
    while (!atomic_load(&shutdown_flag)) {
        struct timespec timeout = {.tv_sec = 0, .tv_nsec = 100000000};
        int signo = sigtimedwait(&sigset, NULL, &timeout);
        if (signo == SIGINT) {
            atomic_store(&shutdown_flag, 1);
            cleanup_shutdown();
            break;
        }
    }
    return NULL;
}
\end{lstlisting}

\subsection{Statistics and Logging}

All operations are logged to stdout with required format. Statistics are maintained using C11 atomic operations for thread safety.

\section{Testing and Validation}

\subsection{Official Test Scenarios}

The implementation was tested against all required scenarios:

\begin{enumerate}
    \item \textbf{Scenario 1}: 10 mixed orders with 3 couriers
    \item \textbf{Scenario 2}: 10 economy orders with 10 couriers  
    \item \textbf{Scenario 3}: SIGINT cancellation with 20 mixed orders
    \item \textbf{Scenario 4}: Memory safety validation with Valgrind
    \item \textbf{Scenario 5}: 20 mixed orders with 10 couriers
\end{enumerate}

\subsection{Additional Stress Tests}

\begin{itemize}
    \item Priority queue ordering verification
    \item SIGINT handling with active deliveries
    \item Invalid input handling
    \item CLI argument validation
    \item Large-scale test (200 orders, 5 couriers)
\end{itemize}

\subsection{Memory Safety}

Valgrind analysis confirmed:
\begin{itemize}
    \item No memory leaks
    \item No invalid memory access
    \item All heap blocks freed
\end{itemize}

\section{Results}

\subsection{Performance Metrics}

Test results show correct behavior across all scenarios:

\begin{table}[H]
\centering
\begin{tabular}{|c|c|c|c|}
\hline
Scenario & Orders & Completed & Cancelled \\
\hline
1 & 10 & 10 & 0 \\
2 & 10 & 10 & 0 \\
3 & 19 & Variable & Variable \\
4 & 10 & 10 & 0 \\
5 & 19 & 19 & 0 \\
\hline
\end{tabular}
\caption{Test Scenario Results}
\end{table}

\subsection{Correctness Verification}

\begin{itemize}
    \item \checkmark Priority queue maintains correct ordering
    \item \checkmark SIGINT properly cancels pending orders
    \item \checkmark Active deliveries complete before shutdown
    \item \checkmark Statistics accurately reflect operations
    \item \checkmark Thread synchronization prevents race conditions
    \item \checkmark Memory management is leak-free
\end{itemize}

\section{Conclusion}

The implementation successfully meets all requirements of CSE344 HW5:

\begin{itemize}
    \item Multi-threaded architecture with proper synchronization
    \item Priority-based order processing
    \item Atomic statistics tracking
    \item Safe SIGINT handling
    \item Required logging format
    \item Memory safety and performance
\end{itemize}

The system demonstrates robust concurrent programming practices and handles edge cases appropriately.

\end{document}